\documentclass[a4paper,12pt]{article}
\usepackage[applemac]{inputenc}
\usepackage[OT1,T1]{fontenc}

\setlength{\textwidth}{160mm}
\setlength{\textheight}{230mm}
\setlength{\oddsidemargin}{-5mm}
\setlength{\topmargin}{-10mm}
% \usepackage{times}
\usepackage{setspace}
\usepackage{lscape}
\usepackage{amsmath,amssymb}
\usepackage{fancyvrb}
\usepackage{multirow}
\usepackage{float}
\usepackage{hyperref}
\usepackage{color}
\usepackage{listings}
\usepackage{fancyhdr}

\usepackage{fancyvrb,moreverb,boxedminipage}

\textwidth=6in
\textheight=8in
\topmargin=0.0in
\oddsidemargin=0in 
\baselineskip=18pt
\headheight=2cm
\setcounter{MaxMatrixCols}{10}

\newcommand{\ns}{\!\!\!}
\newcommand{\e}{\varepsilon}
\newcommand{\be}{\beta}
\newcommand{\s}{\sigma}
\newcommand{\vv}{\vspace{.5cm}}
\DeclareMathOperator{\plim}{plim}
\DefineVerbatimEnvironment{code}{Verbatim}{fontsize=\small}
\DefineVerbatimEnvironment{example}{Verbatim}{fontsize=\small}

\def\by{\mathbf{y}}
\def\bx{\mathbf{x}}
\def\ba{\mathbf{a}}
\def\bb{\mathbf{b}}
\def\be{\mathbf{e}}
\def\bu{\mathbf{u}}
\def\bv{\mathbf{v}}
\def\bz{\mathbf{z}}

%===========
%Greek letter vectors
%===========
\def\bbeta{\mbox{\boldmath $\beta$}}
\def\bgamma{\mbox{\boldmath $\gamma$}}
\def\bvareps{\mbox{\boldmath $\varepsilon$}}
\def\bleta{\mbox{\boldmath $\eta$}}
\def\bmu{\mbox{\boldmath $\mu$}}
\def\btheta{\mbox{\boldmath $\theta$}}
\def\bsigma{\mbox{\boldmath $\sigma$}}
\def\blambgam{\mbox{\boldmath $\lambda_{\gamma}$}}
\def\blambbet{\mbox{\boldmath $\lambda_{\beta}$}}
\def\blambda{\mbox{\boldmath $\lambda$}}

%===========
%Matrices
%===========
\def\bX{\mathbf{X}}
\def\bD{\mathbf{D}(\bsigma)}
\def\bV{\mathbf{V}(\bsigma)}
\def\bVinv{\mathbf{V}^{-1}(\bsigma)}
\def\bR{\mathbf{R}(\bsigma)}
\def\bA{\mathbf{A}}
\def\bB{\mathbf{B}}
\def\bZ{\mathbf{Z}}
\def\bIn{\mathbf{I_{N}}}
\def\bIm{\mathbf{I_{m}}}
\def\bIni{\mathbf{I_{n}}}
\def\SSW{\mathrm{SSW}}
\def\SSB{\mathrm{SSB}}
\def\Normal{\mathcal{N}\!}

%===========
%Parentheses, etc.
%===========
\def\op{\left(}
\def\cp{\right)}
\def\oc{\left[}
\def\cc{\right]}
\def\ok{\left\{}
\def\ck{\right\}}

%===========
%Statistical functions and real numbers.
%===========
\def\exE{\mathbb{E}}
\def\Real{\mathbb{R}}
\def\Var{\mathbb{V}ar}
\def\Proba{\mathbb{P}}

\hypersetup{colorlinks=true,linkcolor=blue,citecolor=blue, urlcolor=blue}
\urlstyle{tt}

\newenvironment{exercices}{\newcounter{moncompteur}\begin{list}
   {\bfseries Ex. \arabic{moncompteur}}{\usecounter{moncompteur}
     \setlength{\labelwidth}{2.6em}\setlength{\leftmargin}{3.1em}}}{\end{list}}
\renewcommand{\labelenumi}{\bfseries\alph{enumi})}
\newcommand{\splus}[3][0]{\texttt{>~#2}\ \ \ \ \textit{#3}\\[#1mm]}

\newenvironment{exo}
{\begin{center}\setlength{\textwidth}{140mm} \underline{Exercise:} \begin{minipage}[t]{120mm}}
{\end{minipage}\setlength{\textwidth}{160mm}\end{center}}


\usepackage{verbatim}

\lstset{ %
  language=R, 
  basicstyle=\ttfamily,
  backgroundcolor=\color{lightgray},   % choose the background color; you must add \usepackage{color} or \usepackage{xcolor}
  xleftmargin= 10 pt,
  xrightmargin= 10 pt,
  escapeinside={\%*}{*)},          % if you want to add LaTeX within your code
  frame=none,                    % adds a frame around the code
  numbers=none,                    % where to put the line-numbers; possible values are (none, left, right)
  numbersep=5pt,                   % how far the line-numbers are from the code
  numberstyle=\tiny\color{mygray}, % the style that is used for the line-numbers
  showspaces=false,                % show spaces everywhere adding particular underscores; it overrides 'showstringspaces'
}

\pagestyle{fancy}

\definecolor{mygreen}{rgb}{0,0.6,0}
\definecolor{mygray}{rgb}{0.5,0.5,0.5}
\definecolor{mymauve}{rgb}{0.58,0,0.82}
\definecolor{lightgray}{gray}{0.95}


%%%%%%%%%%%%%%%%%%%%%%%%%%%%%%%%%%%%%%%%%%%%%%%%%%%%%%%%%%%%%%%%%%%%%%%%%%%%
%
\begin{document}
%

\lhead{University of Geneva\\ \textbf{GLAM}\\
Pr. Eva Cantoni}
\rhead{GSEM\\ Spring 2022 \\ \textbf{Homework 02}}

\begin{center}
\bigskip

\bigskip

\bigskip

\bigskip

\bigskip

\bigskip

\fbox{{\Large Handwritten exercises on GLMs}}

\bigskip

\bigskip

\bigskip
\end{center}

%%%%%%%%%%%%%%%%%%%%%%%%%%%%%%%%%%%%%%%%%%%%%%%%%%%%%%%%%%%%%%%%%%%%%%%%%%%%

\setlength{\parskip}{1ex plus 0.5ex minus 0.2ex}
\parindent=0cm
\setlength{\parskip}{1ex plus 0.5ex minus 0.2ex}
\linespread{1.1}
\sloppy

\section{\underline{Exercise 1}}

Show that the Binomial distribution and the Gamma distribution belong to the exponential family of distributions, by identifying $\theta$, $A$, $\phi$, and the functions $b()$ and $c()$ (cf.\ p.~53 of the course slides). Using the properties of the distributions of the exponential family, give the expectation and the variance in each case.

\textit{Notation:}
\begin{itemize}
\item If $Y$ follows a  binomial distribution ${\mathcal B}(m,p)$, then 
$$P(Y=y; m,p) = {m \choose y} p^{y} (1- p)^{m-y}, \ \  \mbox{ for } y=0, \ldots, m.$$
\item If $Y$ follows a Gamma distrbution $\Gamma(\mu,\nu)$, then 
$$f(y;\mu,\nu)=\frac{\nu \exp(-\nu y/\mu)(\nu y/\mu)^{\nu-1}}{\mu\Gamma(\nu)}, \ \  \mbox{ for } y>0.$$
\end{itemize}


\section{\underline{Exercise 2}}
The beta distribution is used for situations where the measured variable takes values between 0 and 1, for example proportions.

We let $Y$ follow a beta distribution with parameters $\mu$ and  $\phi$, where $\mu$ is the parameter of interest and $\phi$ is a nuisance parameter. The density function of $Y$ is
$$f(y ) = \frac{\Gamma(\phi)}{\Gamma(\phi \mu)\Gamma(\phi (1-\mu) )} y^{\phi \mu - 1} (1-y)^{\phi (1-\mu) -1},  \ \  \mbox{ for } 0 < y < 1.$$

Does this distribution belong to the exponential family? Justify your answer, and if yes deduce $E(Y)$ and $Var(Y)$.

\section{\underline{Exercise 3}}

\begin{itemize}
\item Let $f(y;\theta)$ be the probability density function of a random variable $Y$, for a given parameter $\theta$. Without using the exponential family structure, show that for ${\displaystyle U(\theta;Y) = \frac{\partial}{\partial \theta} \log f(Y;\theta)}$
\begin{enumerate}
    \item{$\exE\left[U(\theta;Y)\right]=0$,}
    \item{${\displaystyle \exE\left[\frac{\partial U(\theta;Y)}{ \partial \theta}\right]=  - \exE\oc U^2(\theta;Y)\cc}$.}
\end{enumerate}
\item What condition do we have to impose on the support of $Y$ to prove these two statements?
\end{itemize}


\section{\underline{Exercise 4}}

In this exercise, we consider the problem of Maximum Likelihood estimation of the parameters of a Poisson random variable with a log link, i.e : 
\begin{equation*}
    g\op\mu_i\cp = \mathrm{ln}\ok \exE \oc Y_i \ | \ \bX_i = \bx_i  \cc\ck = \bx_i^T\bbeta
\end{equation*}

\begin{enumerate}
  \item Write the estimating (score) equations for the proposed model.
  \item Show that 
    \begin{equation*}
   \mathcal J(\bbeta; \by) 
   = \sum_{i=1}^n \exp\{\bx_i^T\beta\} 
   \bx_i \bx_i^T.
   \label{eq:J}
  \end{equation*}
  \item Write the formula for the Newton-Raphson algorithm.  

\end{enumerate}

\end{document}
